\documentclass[10pt,letterpaper]{article}
\usepackage{cv-michaillat}
\usepackage[spanish]{babel}

% Pie de página
\grayfoot{Fabián Belmar \sep Curriculum Vitae}{\thepage}

\begin{document}

% ============================================
% ENCABEZADO
% ============================================
\name{Fabián Belmar}

\begin{center}
\href{mailto:fbelmar@cepchile.cl}{fbelmar@cepchile.cl} \sep
+56 9 76371608 \sep
\href{https://belmarfabian.github.io}{belmarfabian.github.io}\\[3pt]
{\footnotesize
\href{https://orcid.org/0000-0003-4239-1874}{ORCID} \sep
\href{https://www.scopus.com/authid/detail.uri?authorId=57220090371}{Scopus} \sep
\href{https://www.researchgate.net/profile/Fabian-Belmar-3}{ResearchGate} \sep
\href{https://scholar.google.es/citations?user=wr6beBsAAAAJ}{Google Scholar}
}
\end{center}

\vspace{6pt}

% ============================================
% POSICIÓN ACTUAL
% ============================================
\block{Posición Actual}

\entry{2025--}{Investigador Asistente, Centro de Estudios Públicos (CEP), Santiago.}
\hspace{2.5em}Proyecto C22 -- Métodos Digitales. Análisis de procesos sociales mediante ciencia de datos y machine learning.

% ============================================
% FORMACIÓN
% ============================================
\block{Formación Académica}

\entry{2020--2025}{Doctor en Procesos e Instituciones Políticas, Universidad Adolfo Ibáñez.}
\hspace{2.5em}Distinción máxima. Tesis: ``Corrupción al margen del Estado. El caso de la FIFA''.\\
\hspace{2.5em}Director: Dr. Aldo Mascareño.

\vspace{4pt}

\entry{2013--2018}{Licenciado en Ciencia Política y Administración Pública. Administrador Público.}
\hspace{2.5em}Universidad de Talca. Distinción máxima. Mejor egresado de la generación.\\
\hspace{2.5em}Tesis: ``Vinculación no programática en la gestión municipal''. Director: Dr. Mauricio Morales.

% ============================================
% EXPERIENCIA
% ============================================
\block{Experiencia en Investigación}

\entry{2023--2024}{Investigador Pasante, Instituto Gino Germani, Universidad de Buenos Aires.}
\hspace{2.5em}Grupo de Estudios de Culturas Populares Contemporáneas. Directora: Dra. Verónica Moreira.

\vspace{4pt}

\entry{2022--2023}{Asistente de Investigación, FONDECYT Regular \#1220004.}
\hspace{2.5em}``Vinculación entre representantes y representados''. IP: Dr. Mauricio Morales.

\vspace{4pt}

\entry{2018--2023}{Coordinador de Proyectos, Centro de Análisis Político, Universidad de Talca.}
\hspace{2.5em}Investigación, vinculación con el medio, análisis de datos electorales.

\vspace{4pt}

\entry{2018--2020}{Asistente de Investigación, FONDECYT Regular \#1180009.}
\hspace{2.5em}``Los Partidos Demócrata Cristianos en América Latina''. IP: Dr. Mauricio Morales.

% ============================================
% PUBLICACIONES
% ============================================
\block{Artículos en Revistas Indexadas (WoS/Scopus)}

\textit{En prensa / En revisión}

\begin{enumerate}[start=11]
    \pub{Belmar, F., \& Mascareño, A. FIFA: Between couplings and hypertrophy. The evolution of the system in the 20th and 21st centuries. \textit{International Journal of Sport Policy and Politics}. [En revisión, 2da ronda]}
    \pub{Belmar, F., Contreras, G., \& Morales, M. (2025). Who demands clientelism? Differences between individual and organizational demand making. \textit{Political Science Research \& Methods}. [En prensa]}
    \pub{Belmar, F., \& Mascareño, A. (2025). Corruption: An uneven field of research. \textit{Societies}, 15(7), 186.}
    \pub{Belmar, F., Morales, M., \& Villarroel, B. (2025). The Impact of Emotions on Voting Intention in a Constitutional Referendum. \textit{Political Studies Review}. [En prensa]}
\end{enumerate}

\textit{Publicados}

\begin{enumerate}[start=7]
    \pub{Belmar, F., \& Morales, M. (2025). Ethnic solidarity in the non-programmatic distribution of benefits: The case of Mapuche mayors in Chile. \textit{Ethnic and Racial Studies}.}
    \pub{Belmar, F., Lara, B., \& Miguieles, J. (2024). Unpacking electoral competition: Campaign spending, vote distribution and turnout in Chilean mayoral elections. \textit{Policy Studies}.}
    \pub{Belmar, F., Maldonado, F., \& Morales, M. (2024). In search of the missing link: Election law infractions and candidate sanctions. \textit{Public Integrity}.}
    \pub{Belmar, F., Morales, M., \& Villarroel, B. (2023). Writing constitution without parties? The programmatic weakness of party-voter linkages in the Chilean political change. \textit{Politics}, 45(1).}
    \pub{Belmar, F., Contreras, G., Morales, M., \& Troncoso, C. (2023). Demanding non-programmatic distribution. Evidence from local governments in Chile. \textit{Policy Studies}.}
    \pub{Belmar, F., \& Morales, M. (2022). Clientelism, turnout and incumbents' performance in Chilean local government elections. \textit{Social Sciences}, 11(8), 361.}
    \pub{Belmar, F., \& Morales, M. (2020). Clientelismo en la gestión local. El caso de los alcaldes en Chile, 2015-2016. \textit{Política y Sociedad}, 57(2), 567--591.}
\end{enumerate}

% ============================================
% CAPÍTULOS
% ============================================
\block{Capítulos de Libro}

\begin{enumerate}
    \pub{Belmar, F. (2024). Descentralización y participación ciudadana en Chile: evaluación de la autonomía de los gobiernos locales. En VV.AA. \textit{La Revolución Silente: Descentralización en Chile}.}
    \pub{Belmar, F. (2020). Las tecnologías al servicio de la comunicación y el funcionamiento interno de los parlamentos. En VV.AA. \textit{Parlamentos Conectados}. Universidad de Oviedo, España.}
    \pub{Belmar, F. (2020). La estrategia de comunicación de los representantes: la construcción de la imagen de los políticos. En VV.AA. \textit{Parlamentos Conectados}. Universidad de Oviedo, España.}
    \pub{Belmar, F., Herrera, M., \& Obregón, M. (2018). El efecto de la elección de concejales sobre la elección de consejeros regionales 2017. En Morales, M., Navarrete, B. \& Vial, C. \textit{Política Subnacional en Chile}. RIL Editores.}
\end{enumerate}

% ============================================
% DOCENCIA
% ============================================
\block{Docencia}

\entry{2025}{Universidad de Talca}
\hspace{2.5em}Evaluación de Políticas, Programas y Proyectos; Ciencia de la Administración Pública.

\vspace{2pt}

\entry{2025}{Universidad Diego Portales}
\hspace{2.5em}Métodos Cuantitativos Introductorios; Métodos Cuantitativos Avanzados.

\vspace{2pt}

\entry{2025}{INACAP}
\hspace{2.5em}Introducción a las Políticas Públicas; Sistemas de Compras Públicas.

\vspace{2pt}

\entry{2023--2024}{Universidad Central}
\hspace{2.5em}Análisis de Datos Cuantitativos; Métodos Cuantitativos para la Investigación.

\vspace{2pt}

\entry{2021--2022}{Universidad de Talca}
\hspace{2.5em}Métodos Cualitativos para las Ciencias Sociales; Análisis de Datos para las Ciencias Sociales.

\vspace{2pt}

\entry{2022}{Universidad Adolfo Ibáñez}
\hspace{2.5em}Redes Sociales, Democracia y Ciudadanía Digital.

% ============================================
% PONENCIAS
% ============================================
\block{Ponencias (Selección)}

\begin{itemize}
    \item (2024) I Congreso Internacional de Evaluación de la Gobernanza. Universidad Católica de Temuco, Chile.
    \item (2024) I Congreso RELASSC. Recife, Brasil.
    \item (2023) XIV Congreso Chileno de Ciencia Política. Universidad de Santiago, Chile.
    \item (2023) VIII Congreso Uruguayo de Ciencia Política. Universidad Católica del Uruguay.
    \item (2022) IV Congreso Internacional ICCA. Universidad Complutense de Madrid, España.
    \item (2022) V Congreso Latinoamericano FLACSO. Montevideo, Uruguay.
\end{itemize}

% ============================================
% BECAS
% ============================================
\block{Becas, Distinciones y Fondos}

\begin{itemize}
    \item \textbf{2023--2024}: Beca Pasantía Internacional ANID (12,000 USD). IIGG, Universidad de Buenos Aires.
    \item \textbf{2021--2022}: Beca Beneficios Complementarios ANID (6,000 USD).
    \item \textbf{2020}: Beca Doctorado Nacional ANID (8,000 USD).
    \item \textbf{2019}: Fondo de Desarrollo Institucional, Ministerio de Educación (7,000 USD).
    \item \textbf{2018}: Mejor Egresado de la Generación, Universidad de Talca.
\end{itemize}

% ============================================
% COMPETENCIAS
% ============================================
\block{Software e Idiomas}

\noindent\textbf{Programación:} R (Avanzado) \sep Python (Intermedio) \sep SQL (Intermedio) \sep \LaTeX\ (Avanzado)

\noindent\textbf{Análisis cuantitativo:} SPSS \sep STATA \sep R \sep Jamovi (Avanzado)

\noindent\textbf{Análisis geoespacial:} QGIS (Avanzado)

\noindent\textbf{Análisis cualitativo:} ATLAS.ti (Avanzado) \sep VOSviewer (Intermedio)

\noindent\textbf{Idiomas:} Español (Nativo) \sep Inglés (Intermedio escrito, Básico oral)

% ============================================
% ÁRBITRO
% ============================================
\block{Árbitro en Revistas Académicas}

\textit{Ethnic and Racial Studies} \sep \textit{Journal of Infrastructure, Policy and Development} \sep \textit{PLOSone} \sep \textit{Economía y Política} \sep \textit{Economies} \sep \textit{Sustainability} \sep \textit{Revista Desarrollo y Sociedad} \sep \textit{Revista Iberoamericana de Estudios Municipales (RIEM)}

% ============================================
% REFERENCIAS
% ============================================
\block{Referencias}

\noindent\textbf{Mauricio Morales}. Universidad de Talca. Doctor en Ciencia Política, PUC. \href{mailto:mmoralesq@utalca.cl}{mmoralesq@utalca.cl}

\noindent\textbf{Aldo Mascareño}. Universidad Adolfo Ibáñez. Doctor en Sociología, U. Bielefeld. \href{mailto:amascareno@cepchile.cl}{amascareno@cepchile.cl}

\noindent\textbf{Bernardita Escobar}. Universidad de Valparaíso. Doctora en Economía, Cambridge. \href{mailto:beni.escobar.andrae@gmail.com}{beni.escobar.andrae@gmail.com}

\noindent\textbf{Daniel Chernilo}. Universidad Adolfo Ibáñez. Doctor en Sociología, U. Warwick. \href{mailto:daniel.chernilo.s@uai.cl}{daniel.chernilo.s@uai.cl}

\end{document}
